\documentclass[12pt,a4paper]{report}
\usepackage[backend=bibtex,style=numeric,sorting=none]{biblatex}
\usepackage[utf8]{inputenc}
\usepackage[T1]{fontenc}
\usepackage[ngerman]{babel}
\usepackage{amsmath, amssymb, amsthm} % Mathe
\usepackage{geometry}
\geometry{a4paper, left=3cm, right=2.5cm, top=2.5cm, bottom=3cm}
\usepackage{graphicx}
\usepackage{hyperref}
\usepackage{csquotes}
\usepackage{biblatex}
\usepackage{subcaption}
\usepackage{listings}
\usepackage{pdfpages}
\usepackage[absolute,overlay]{textpos}
\usepackage{hyperref}
\addbibresource{literatur.bib}

\newtheorem{theorem}{Satz}[chapter]


\newtheorem{definition}[theorem]{Definition}
\newtheorem{lemma}[theorem]{Lemma}
\newtheorem{korollar}[theorem]{Korollar}
\newtheorem{bemerkung}[theorem]{Bemerkung}
\newtheorem{satz}[theorem]{Satz}

\theoremstyle{remark}
\newtheorem{beispiel}[theorem]{Beispiel}



\begin{document}


\vspace*{2cm}
\chapter*{Ist die Challange vollständig?}

Die Challange enthält alles was ich brauche um den Angriff durchzuführen. Ich habe den Public-Key, U, V und eine 
funktionierende Decryption-Funktion.

\chapter{Ist es möglich, den geheimen Schlüssel aus dem öffentlichen Schlüssel zu extrahieren?}

Ich konnte den geheimen Schlüssel extrahieren. Ich habe den Schlüssel vier mal extrahiert und brauchte dabei im 
Schnitt ca. 6000 Iterationen. (Eine Iteration ist das bestimmen einer Permuationsmatrix H für H * P).

\chapter{Ist es klar, wie die Entschlüsselungs-Funktion aufgerufen werden muss?}

Ja. 

\chapter{Gibt die Entschlüsselungsfunktion bei Eingabe des Chiffretexts und des geheimen Schlüssels den Klartext wieder?}

Ja. Er lautet: 

    
\chapter{Kann irgendetwas verbessert werden?}

\begin{itemize}
    \item Im Code kann mehr mit Variablen gearbeitet werden. Die Decryption-Funktion ist recht unübersichtlich, da 
    U und V nicht als Variablen in die Funktion eingefügt werden, die wurden nur reinkoppiert, was sehr unübersichtlich 
    ist.
    \item Sk war noch in der Decryption-Funktion. Sollte entfernt werden.
    \item Der Public-Key war s,h und nicht h, s.
\end{itemize}


\end{document}

